%% raw_s3b.tex -- raw transcription of Slide Set 3, PDF pages 43--93
%% (Discriminant analysis: LDA, QDA, comparison of classifiers). Trimmed + compressed: only material not in Part I.
\section{Slide Set 3 (cont.) --- Discriminant analysis: LDA \& QDA (pp.~43--93; past the Slide-42 cutoff)}
\ixh{discriminant analysis}\ixh{linear discriminant analysis (LDA)}\ixh{quadratic discriminant analysis (QDA)}
{\color{figgray}\itshape [Omitted (in Part I): S3.43--44, 46, 48--51, 53, 56, 58--60, 62, 65--74, 76--79, 81--84, 86--93.
Past the Slide-42 cutoff; kept because PE Q20, Lab~3, HW~3 test LDA/QDA.]}

\subsection{Why LDA?}
\ixh{linear discriminant analysis (LDA)}

\slide{3}{45}{Why bother considering LDA?}
\ixh{perfect separation}
\begin{routput}
> x <- seq(100)
> y <- c(rep(0, 50), rep(1, 50))
> fit <- glm(y~x, family = "binomial")        (*@\ixch{glm()}\ixch{family = "binomial"}@*)
\end{routput}

\slide{3}{47}{[feature extraction vs feature selection]}
\fig{Left, ``\ix{feature extraction}'': original features transformed/combined into a few \emph{new} features (rows of
5 $\to$ 8 $\to$ 6 $\to$ 3 shapes). Right, ``\ix{feature selection}'': row of 7 circles, 3 crossed out --- a subset of the
\emph{original} features kept.}\ixh{dimension reduction}
{\color{figgray}\itshape [note: LDA = feature extraction --- new axes LD1, LD2, \ldots{} from combinations of predictors; S3.64--67, 80--85.]}

\subsection{Geometric picture: gene-expression example}
\ixh{projection}\ixh{gene-expression example}

\slide{3}{52--57}{LDA: A motivating example}
\fig{Number line ``Transcription for Gene X'' (fewer $\leftarrow\to$ more transcripts), 12 dots left to right:
G G G G R G R G R R R R (green = drug works, red = does not). \url{https://youtu.be/azXCzI57Yfc}}
\begin{itemize}
\item S3.54: Mostly, drug works with low Gene X transcription (4 leftmost green circled), not with high (4 rightmost red).
\item S3.55: But overlap\ixh{class overlap} (middle R G R G) and no obvious ``cutoff'' for who gets the drug.
\item S3.57: Gene X does an OK job. Can we do better with more than one gene?
\end{itemize}

\slide{3}{61--64}{LDA: A motivating example}
\begin{itemize}
\item S3.61: 2D \ix{scatterplot} Gene X vs.\ Gene Y: 6 green mostly lower-left, 6 red along top/upper-right (+1 red lower
right); a diagonal dashed line (linear \ix{decision boundary}) separates them. Two genes separate better. Three?
\item S3.63: With 3 genes (X, Y, Z) we use a \ix{plane}\ixh{hyperplane} to separate the categories.
\item S3.64: Reducing 2-D graph $\to$ 1-D graph (number line)\ixh{dimension reduction}: G G G G G R R G R R R R.
\end{itemize}

\slide{3}{75}{[distance between means alone vs.\ means and scatter]}
\ixh{distance between means}\ixh{scatter}
\fig{2D: 6 red in an upper band, 6 green in a lower band, both rising gently right. Projected onto a line along the
bands: means far apart but classes interleave in the middle. Projected onto a near-vertical line: red tight at top,
green tight at bottom, no overlap.}
Maximizing only the distance between means $\to$ poor separation; optimizing distance between means \emph{and}
scatter $\to$ nice separation.

\slide{3}{80}{[LDA with three genes]}
\ixh{LD1}\ixh{linear discriminant (LD1, LD2)}
\fig{3D Gene X/Y/Z cloud with a diagonal line = the new axis LDA created; data projected onto it: 6 green
lower-left, 6 red upper-right, completely separated.}

\slide{3}{85}{[LDA plot: LD1 vs LD2]}
\ixh{linear discriminant (LD1, LD2)}
\fig{R plot ``LDA'', \ix{LD1} ($-4$ to 2) vs.\ \ix{LD2} ($-2$ to 2), three classes: red at left (LD1 $-4.2$ to $-0.8$,
LD2 $-1$ to 1); black upper right (LD1 0--2, LD2 $-0.7$ to 2.6); blue lower right (LD1 $-0.2$ to 2.2, LD2 $-2.7$ to 0),
a few blue among black. LD1 separates red from the rest; LD2 separates black from blue.}
